% Pista geo-23s-q5 · Geometria
% Procedència: PAU Catalunya 2023 · Convocatòria de setembre · Sèrie 2 · Qüestió 5
\documentclass[11pt,a4paper]{article}
\usepackage{pista}
\pistainfo{Catalunya 2023}{Convocatòria de setembre}{Sèrie 2}

\begin{document}

\section*{\textcolor{blauPista}{Qüestió 5}}

\noindent Siguin $r_{1}$ i $r_{2}$ les rectes definides per $r_{1}: x - 1 = y = -z$ i per $r_{2}: x = y = z$, respectivament.

\subsection*{\textit{a)} \normalfont Calculeu l'equació paramètrica de la recta que talla perpendicularment les rectes $r_{1}$ i $r_{2}$. \hspace*{\fill}\textnormal{[1,75~punts]}}

\pista{Comença identificant un vector director i un punt de cada recta a partir de les equacions contínues.

\medskip
La perpendicular comuna queda determinada per dues coses:
\begin{itemize}
\item \textit{Direcció.} Ha de ser perpendicular alhora a $\vec{v}_{1}$ i a $\vec{v}_{2}$: és $\vec{v}_{1} \times \vec{v}_{2}$.
\item \textit{Un punt.} Pren un punt genèric $P_{1}$ d'$r_{1}$ (paràmetre $t$) i un punt genèric $P_{2}$ d'$r_{2}$ (un paràmetre diferent, $s$), i imposa que $\overrightarrow{P_{1} P_{2}}$ sigui perpendicular alhora a $\vec{v}_{1}$ i a $\vec{v}_{2}$:
\[
\overrightarrow{P_{1} P_{2}} \cdot \vec{v}_{1} = 0, \qquad \overrightarrow{P_{1} P_{2}} \cdot \vec{v}_{2} = 0
\]
Resol aquest sistema de dues equacions amb dues incògnites per obtenir $t$ i $s$, i per tant $P_{1}$ i $P_{2}$.
\end{itemize}
La recta que uneix $P_{1}$ i $P_{2}$ és la perpendicular comuna.}

\subsection*{\textit{b)} \normalfont Calculeu la distància entre $r_{1}$ i $r_{2}$. \hspace*{\fill}\textnormal{[0,75~punts]}}

\pista{La distància entre les dues rectes és la longitud del segment de perpendicular comuna que va de $P_{1}$ a $P_{2}$:
\[
d(r_{1}, r_{2}) = \left|\overrightarrow{P_{1} P_{2}}\right|
\]
Els punts $P_{1}$ i $P_{2}$ ja els has trobat a l'apartat anterior: només et cal calcular aquest mòdul.}

\end{document}
